\documentclass[11pt,a4paper]{article}
\usepackage[utf8]{inputenc}
\usepackage[italian]{babel}
\usepackage{geometry}
\geometry{margin=2cm}
\usepackage{hyperref}
\usepackage{enumitem}
\usepackage{xcolor}
\usepackage{titlesec}

\definecolor{primary}{HTML}{2E3440}
\definecolor{accent}{HTML}{5E81AC}

\hypersetup{
    colorlinks=true,
    linkcolor=accent,
    urlcolor=accent,
}

\titleformat{\section}{\Large\bfseries\color{primary}}{}{0em}{}[\titlerule]
\titlespacing{\section}{0pt}{12pt}{6pt}

\begin{document}

% Header
\begin{center}
    {\LARGE\bfseries\color{primary} Riccardo Figliozzi}\\[4pt]
    {\large AI Trainer | AI Consultant}\\[4pt]
    Firenze, Toscana, Italia
\end{center}

\vspace{4pt}

% Contatti
\begin{center}
    \href{mailto:riccardo.figliozzi@gmail.com}{riccardo.figliozzi@gmail.com} \quad
    \href{https://linkedin.com/in/riccardofigliozzi}{LinkedIn} \quad
    \href{https://riccardofigliozzi.com}{Sito Personale}
\end{center}

\vspace{4pt}

% Competenze Principali
\section{Competenze Principali}
\begin{itemize}[leftmargin=*, nosep]
    \item \textbf{Strumenti:} n8n, Generative AI, Data Analytics
    \item \textbf{Linguaggi:} Python, R, SQL
    \item \textbf{AI/ML:} LangChain, LangGraph, LLM, RAG, Prompt Engineering
\end{itemize}

% Lingue
\section{Lingue}
\begin{itemize}[leftmargin=*, nosep]
    \item \textbf{Italiano:} Madrelingua
    \item \textbf{Inglese:} Professionale completo
    \item \textbf{Francese:} Professionale completo
    \item \textbf{Spagnolo:} Professionale completo
\end{itemize}

% Formazione
\section{Formazione}
\begin{itemize}[leftmargin=*, nosep]
    \item \textbf{Universit\`{a} Paris Dauphine - PSL} \hfill Settembre 2020 -- 2021\\
    Master 2 (M2), MIAGE Informatique D\'{e}cisionnelle

    \item \textbf{Universit\`{a} di Pisa} \hfill Settembre 2019 -- Settembre 2021\\
    Magistrale in Data Science and Business Informatics

    \item \textbf{Universit\`{a} degli Studi di Messina} \hfill Settembre 2018 -- Novembre 2019\\
    Master in Banking and Finance, Economia

    \item \textbf{Universidad de Murcia} \hfill Settembre 2017 -- Luglio 2018\\
    Marketing e comunicazione

    \item \textbf{Universit\`{a} degli Studi di Messina} \hfill Settembre 2015 -- Ottobre 2018\\
    Scienze dell'informazione: comunicazione pubblica e tecniche giornalistiche
\end{itemize}

% Esperienza Professionale
\section{Esperienza Professionale}

\textbf{Data Masters} \hfill Febbraio 2025 -- Presente\\
\textit{AI Teacher / AI Engineer}
\begin{itemize}[leftmargin=*, nosep]
    \item Progettazione e sviluppo di chatbot agentic per l'onboarding utenti sulla piattaforma Data Masters, utilizzando LangChain, LangGraph, LangSmith e ChromaDB (architettura RAG).
    \item Workflow di automazione AI basati su LLM e agenti usando n8n: pipeline di web scraping, elaborazione dati, reportistica automatizzata; agente RAG per il customer service.
    \item Insegnamento di AI generativa, data science e machine learning: trasferimento di competenze a professionisti e aziende.
\end{itemize}

\vspace{4pt}

\textbf{Freelance} \hfill Agosto 2024 -- Presente\\
\textit{AI Adoption Consultant \& Trainer} -- Firenze, Toscana, Italia
\begin{itemize}[leftmargin=*, nosep]
    \item Analisi dei processi aziendali per identificare opportunit\`{a} di automazione tramite soluzioni di Intelligenza Artificiale.
    \item Progettazione e sviluppo di soluzioni AI (agenti, workflow di automazione, integrazioni con strumenti low-code/no-code e piattaforme AI).
    \item Utilizzo e configurazione di strumenti e framework chiave per AI generativa e automazione (modelli LLM, piattaforme di orchestrazione, RPA, strumenti di prompt engineering).
    \item Supporto al team nella preparazione di proposte di progetto e PoC (Proof of Concept) per i clienti.
    \item Guida del client nell'adozione delle nuove soluzioni e degli strumenti AI introdotti.
    \item Monitoraggio continuo dell'evoluzione degli strumenti e delle tecnologie AI per proporre soluzioni aggiornate ed efficaci.
    \item Preparazione di documentazione tecnica e report di avanzamento progetto.
\end{itemize}

\vspace{4pt}

\textbf{Accademia del Levante} \hfill Febbraio 2025 -- Marzo 2025\\
\textit{Docente}
\begin{itemize}[leftmargin=*, nosep]
    \item Contributo alla formazione di futuri Data Engineer insegnando i moduli di Excel e Statistica in un corso online specializzato.
\end{itemize}

\vspace{4pt}

\textbf{Beyond Knowledge} \hfill Aprile 2023 -- Luglio 2024\\
\textit{Data Scientist} -- Firenze
\begin{itemize}[leftmargin=*, nosep]
    \item Sviluppo di soluzioni di machine learning classiche e avanzate, collaborazione con data scientist per identificare problemi, selezionare algoritmi, preparare dati, addestrare/valutare modelli e deployarli in produzione.
    \item Manutenzione di sistemi back-end utilizzando Django, RestAPI, n8n e Docker, garantendo un'integrazione fluida tra modelli ML e sistemi esistenti.
\end{itemize}

\vspace{4pt}

\textbf{SDG Group Italy} \hfill Febbraio 2022 -- Aprile 2023\\
\textit{Data Engineer} -- Firenze, Toscana, Italia
\begin{itemize}[leftmargin=*, nosep]
    \item Costruzione e analisi di flussi dati utilizzando pipeline ETL e migrazione senza problemi su piattaforme cloud come AWS.
\end{itemize}

\vspace{4pt}

\textbf{OR2S -- Observatoire r\'{e}gional de la sant\'{e} et du social} \hfill Aprile 2021 -- Ottobre 2021\\
\textit{Data Scientist} -- Rouen, Normandia, Francia
\begin{itemize}[leftmargin=*, nosep]
    \item Sviluppo di un processo automatizzato utilizzando tecniche di machine learning non supervisionato in Python, per estrarre insights significativi e pattern da dati non etichettati, avvalendosi sia di tecniche di clustering che di riduzione della dimensionalit\`{a}.
\end{itemize}

\vspace{4pt}

\textbf{ARDOXSO} \hfill Gennaio 2020 -- Giugno 2020\\
\textit{Search Engine Optimization Copywriter} -- Pisa, Toscana, Italia
\begin{itemize}[leftmargin=*, nosep]
    \item Produzione di post sul blog ottimizzati SEO su vari temi di marketing digitale.
    \item Analisi dei dati per identificare argomenti di tendenza e parole chiave, informando la strategia dei contenuti.
    \item Utilizzo di strumenti di web analytics per monitorare le prestazioni dei contenuti e migliorare iterativamente le strategie SEO.
\end{itemize}

\vspace{4pt}

\textbf{Studio Legale Cavalletti} \hfill Novembre 2019 -- Aprile 2020\\
\textit{Social Media Manager} -- Pisa, Toscana, Italia
\begin{itemize}[leftmargin=*, nosep]
    \item Gestione della presenza sui social media di un professionista legale locale, supervisione e sviluppo di contenuti per LinkedIn e Facebook.
    \item Analisi delle metriche per ottimizzare le prestazioni dei post e la copertura.
\end{itemize}

\vspace{4pt}

\textbf{Idib Group} \hfill Novembre 2018 -- Giugno 2019\\
\textit{Data Analyst} -- Messina, Sicilia, Italia
\begin{itemize}[leftmargin=*, nosep]
    \item Lavoro con big data utilizzando strumenti come Figma, Miro e inizio della programmazione con R e Python per Business Intelligence e analisi di mercato.
\end{itemize}

\vspace{4pt}

\textbf{UniVersoMe} \hfill Dicembre 2016 -- Luglio 2017\\
\textit{Social Media Manager} -- Messina, Sicilia, Italia
\begin{itemize}[leftmargin=*, nosep]
    \item Social Media Manager di questo progetto universitario realizzato da studenti. Esperienza che ha permesso la crescita professionale e personale, migliorando la creativit\`{a} e molte altre soft-skill.
\end{itemize}

% Certificazioni
\section{Certificazioni}
\begin{itemize}[leftmargin=*, nosep]
    \item Google Analytics Individual Qualification
    \item DELE B2
    \item Rome Business Game
    \item English C1
\end{itemize}

\end{document}
